\chapter{Terminal Obstructions}

\section{Purpose of Terminal Obstructions}

A terminal obstruction is a configuration pattern that prevents an admissible
refinement process from reaching a terminal state too quickly. It is not merely
a hard example. It is a structured witness showing that the hypotheses of the
framework produce an actual obstruction to collapse.

In URF language, a terminal obstruction supplies three pieces of data:

\begin{enumerate}
\item a family of configurations with persistent unresolved entropy;
\item an admissible refinement class with bounded capacity;
\item a terminal condition that cannot be reached without eliminating the
unresolved entropy.
\end{enumerate}

The obstruction is terminal because it is measured against the endpoint of the
refinement process. A configuration is not obstructive merely because it looks
complicated. It is obstructive when its unresolved structure must be removed
before the terminal condition can hold.

\section{Obstruction Data}

\begin{definition}[Terminal obstruction datum]
A terminal obstruction datum is a tuple
\[
(X,\mathcal R,H,C,T,\mathcal O),
\]
where \(X\) is a state space, \(\mathcal R\) is an admissible refinement class,
\(H\) is an entropy functional, \(C\) is a capacity bound, \(T\subseteq X\) is
a terminal set, and \(\mathcal O\subseteq X\) is a distinguished obstruction
class.
\end{definition}

The obstruction class \(\mathcal O\) is useful only if membership in
\(\mathcal O\) forces a lower bound on terminal refinement depth.

\begin{definition}[Depth-forcing obstruction]
An obstruction class \(\mathcal O\) is depth-forcing with rate \(g(n)\) if
there is a family \(x_n\in\mathcal O\) such that every admissible trajectory
from \(x_n\) to \(T\) has length at least \(g(n)\).
\end{definition}

Depth-forcing is always relative to the admissible refinement class. A
configuration may obstruct local refinement while being easy for a nonlocal
process.

\section{Cyclone Obstruction}

The Cyclone obstruction is a model class for cyclically distributed
indistinguishability. Its role is to represent configurations in which local
views repeat around a global cycle, while the global structure remains
unresolved.

\begin{definition}[Cyclone configuration, schematic form]
A Cyclone configuration of length \(n\) is a configuration
\[
x=(x_0,x_1,\ldots,x_{n-1})
\]
with cyclic indexing such that admissible local observations at bounded radius
cannot distinguish many cyclic shifts or phase placements of the configuration.
\end{definition}

The word ``Cyclone'' is a name for the obstruction pattern, not a claim that
all cycle-like objects have the same rigidity behavior. A concrete theorem
must specify the state space, local observables, entropy, and refinement rule.

\begin{definition}[Cyclone entropy]
For a Cyclone family whose unresolved cyclic alternatives form a set
\(\Omega_n\), the Cyclone entropy is
\[
H_{\mathrm{cyc}}(n)=\log |\Omega_n|.
\]
\end{definition}

If \(|\Omega_n|\) grows exponentially in \(n\), then the Cyclone entropy grows
linearly in \(n\). If it grows polynomially, the entropy grows logarithmically.
The obstruction strength depends on this growth.

\section{Cyclone Depth Bound}

\begin{theorem}[Cyclone Depth Bound]
Let \(x_n\) be a Cyclone family with unresolved alternative set \(\Omega_n\).
Suppose terminal resolution requires distinguishing the alternatives in
\(\Omega_n\), and suppose every admissible refinement step removes at most
\(C>0\) units of Cyclone entropy. Then
\[
\ED(x_n)\ge \frac{\log |\Omega_n|}{C}.
\]
\end{theorem}

\begin{proof}
The unresolved entropy of the Cyclone family is
\[
H_{\mathrm{cyc}}(n)=\log |\Omega_n|.
\]
By hypothesis, each admissible step removes at most \(C\) units of this
entropy. Terminal resolution requires eliminating all unresolved alternatives.
The EntropyDepth lower bound therefore gives
\[
\ED(x_n)\ge \frac{H_{\mathrm{cyc}}(n)}{C}
=
\frac{\log |\Omega_n|}{C}.
\]
\end{proof}

\begin{corollary}[Linear Cyclone obstruction]
If \(|\Omega_n|\ge b^n\) for some \(b>1\), then
\[
\ED(x_n)=\Omega(n).
\]
\end{corollary}

\begin{proof}
The assumption gives
\[
\log |\Omega_n|\ge n\log b.
\]
Applying the Cyclone Depth Bound yields
\[
\ED(x_n)\ge \frac{n\log b}{C}.
\]
\end{proof}

\section{Rigidity Wall}

A rigidity wall is a terminal obstruction that cannot be crossed by bounded
refinement without paying depth.

\begin{definition}[Rigidity wall]
A rigidity wall is a family \(x_n\in X_n\) such that:
\begin{enumerate}
\item \(H(x_n)\) grows with \(n\);
\item every admissible step removes at most \(C\) entropy;
\item terminal resolution requires eliminating the entropy measured by \(H\).
\end{enumerate}
\end{definition}

\begin{theorem}[Terminal Rigidity Wall]
Let \(x_n\) be a rigidity-wall family with
\[
H(x_n)\ge g(n).
\]
If every admissible refinement step removes at most \(C>0\) entropy and the
terminal set requires \(H=0\), then
\[
\ED(x_n)\ge \frac{g(n)}{C}.
\]
\end{theorem}

\begin{proof}
This is the EntropyDepth lower bound applied to the family \(x_n\). Since
\(H(x_n)\ge g(n)\), every terminal trajectory has length at least
\(H(x_n)/C\), and hence at least \(g(n)/C\).
\end{proof}

The theorem is deliberately abstract. It becomes a substantive theorem only
after \(H\), \(C\), \(T\), and the admissible refinement class are instantiated.

\section{Graph Obstruction Example}

Let \(G_n\) be a bounded-degree graph family. Suppose a bounded-radius local
refinement process cannot distinguish among a family of \(N_n\) globally
different placements, labelings, or cyclic phases.

Define
\[
H(G_n)=\log N_n.
\]
If each admissible local refinement round removes at most \(C\) units of this
entropy, then terminal separation requires at least
\[
\frac{\log N_n}{C}
\]
rounds.

\begin{example}[Cycle-phase obstruction]
Consider a graph construction in which a local pattern is repeated around a
cycle and the global phase is not visible to bounded-radius neighborhoods.
If there are \(n\) possible phases, the unresolved phase entropy is
\[
\log n.
\]
If there are exponentially many compatible phase assignments, the entropy is
linear in \(n\).
\end{example}

This example is not a graph-isomorphism lower bound. It is a template showing
what data a graph obstruction must provide before a lower bound can be claimed.

\section{Hypergraph Obstruction Example}

Hypergraphs can support obstruction patterns in which local incidences are
uniform while global incidence structure remains unresolved.

\begin{definition}[Incidence-local obstruction]
An incidence-local obstruction is a hypergraph family for which bounded
incidence neighborhoods do not determine the global hyperedge arrangement
needed for terminal resolution.
\end{definition}

If the number of globally compatible hyperedge arrangements is \(M_n\), then
the unresolved entropy is
\[
H_n=\log M_n.
\]
A bounded incidence-local refinement process with capacity \(C\) then has
terminal depth at least
\[
\frac{\log M_n}{C}.
\]

\section{Spectral Obstruction Example}

A spectral terminal obstruction occurs when unresolved components remain in
modes that the admissible process cannot separate at sufficient rate.

\begin{definition}[Spectral obstruction]
Let \(L_n\) be a family of nonnegative operators and let \(U_n\) be a subspace
of unresolved modes. A spectral obstruction occurs when terminal resolution
requires eliminating or separating \(U_n\), while admissible spectral
refinement removes only bounded information per step.
\end{definition}

If the effective unresolved rank is \(r_n\), one may use
\[
H_n=\log r_n
\]
as a rank entropy. If the refinement capacity is \(C\), the terminal depth is
bounded below by \(H_n/C\), provided the rank entropy is the correct terminal
quantity.

\section{Obstruction Versus Example}

Not every example is an obstruction. An example becomes an obstruction only
when it proves a lower bound against a declared refinement class.

\begin{center}
\begin{tabular}{ll}
\toprule
Object & Required evidence \\
\midrule
Example & illustrates a pattern \\
Candidate obstruction & suggests unresolved entropy \\
Certified obstruction & proves entropy growth and capacity bound \\
Terminal obstruction & proves terminal depth lower bound \\
\bottomrule
\end{tabular}
\end{center}

This distinction prevents overclaiming. A named construction should not be
called terminal unless the terminal lower bound has actually been derived.

\section{Audit Checklist}

A terminal-obstruction claim should supply:

\begin{center}
\begin{tabular}{ll}
\toprule
Field & Required content \\
\midrule
Obstruction class & explicit family \(\mathcal O\) or \(x_n\) \\
Local view & admissible observations \\
Unresolved alternatives & set \(\Omega_n\) or equivalent object \\
Entropy growth & lower bound on \(\log |\Omega_n|\) or \(H(x_n)\) \\
Capacity bound & one-step refinement limit \(C\) \\
Terminal condition & required endpoint criterion \\
Depth conclusion & resulting lower bound on \(\ED(x_n)\) \\
Boundary & statement of what is not proved \\
\bottomrule
\end{tabular}
\end{center}

\section{Boundary}

This chapter defines terminal obstruction templates and proves only the
capacity-depth lower bounds that follow from explicit entropy and capacity
hypotheses. It does not prove that Cyclone configurations exist in every
domain, does not prove graph or hypergraph lower bounds without the required
instantiation data, and does not prove unrestricted Chronos-RR, unrestricted
H4.1/FGL, P vs NP, or any Clay problem.
